\documentclass[10pt,twocolumn]{article}
\usepackage[margin=0.6in]{geometry}
\usepackage{amsmath, amssymb, amsthm, enumitem, booktabs, array}

\setlength{\parindent}{0pt}
\setlength{\parskip}{4pt}
\pagestyle{empty}

\newcommand{\sect}[1]{\medskip\noindent\textbf{\large #1}\medskip\hrule\smallskip}
\newcommand{\subsect}[1]{\smallskip\noindent\textbf{#1}\par}

\begin{document}

\begin{center}
{\LARGE\bfseries GRE Math Subject Test}\\[2pt]
{\Large Abstract Algebra Concept Sheet}
\end{center}

\vspace{-4pt}

%======================================================================
\sect{1. Groups --- Definitions}
%======================================================================

\subsect{Group Axioms}
A \textbf{group} $(G,\cdot)$ satisfies:
\begin{itemize}[nosep,leftmargin=*]
\item \textbf{Closure:} $a,b\in G \Rightarrow ab\in G$.
\item \textbf{Associativity:} $(ab)c = a(bc)$.
\item \textbf{Identity:} $\exists\, e\in G$ s.t.\ $ea=ae=a$ for all $a$.
\item \textbf{Inverse:} $\forall\, a\in G,\;\exists\, a^{-1}$ s.t.\ $aa^{-1}=a^{-1}a=e$.
\end{itemize}
If $ab=ba$ for all $a,b$, the group is \textbf{abelian}.

\subsect{Common Examples}
\begin{itemize}[nosep,leftmargin=*]
\item $(\mathbb{Z},+)$, $(\mathbb{Q}^*,\times)$, $(\mathbb{R}^*,\times)$: abelian.
\item $(\mathbb{Z}/n\mathbb{Z},+)$: cyclic group of order $n$.
\item $(S_n, \circ)$: symmetric group, $|S_n|=n!$, non-abelian for $n\ge 3$.
\item $\text{GL}_n(\mathbb{R})$: invertible $n\times n$ matrices, non-abelian for $n\ge 2$.
\item $D_n$: dihedral group of order $2n$ (symmetries of regular $n$-gon).
\item Klein four-group $V_4 = \mathbb{Z}/2\mathbb{Z}\times\mathbb{Z}/2\mathbb{Z}$: abelian, order 4, not cyclic.
\end{itemize}

\subsect{Order of an Element}
$\operatorname{ord}(a) = $ smallest positive $n$ s.t.\ $a^n = e$.\\
If no such $n$ exists, $a$ has \textbf{infinite order}.
\begin{itemize}[nosep,leftmargin=*]
\item $\operatorname{ord}(a)$ divides $|G|$ (by Lagrange).
\item $a^k=e \iff \operatorname{ord}(a)\mid k$.
\item $\operatorname{ord}(a^k)=\dfrac{\operatorname{ord}(a)}{\gcd(k,\operatorname{ord}(a))}$.
\end{itemize}

%======================================================================
\sect{2. Subgroups}
%======================================================================

\subsect{Subgroup Test}
$H\subseteq G$ is a subgroup if:
\begin{itemize}[nosep,leftmargin=*]
\item $H\neq\emptyset$, and
\item $a,b\in H \Rightarrow ab^{-1}\in H$ \quad (one-step test).
\end{itemize}
For finite $H\subseteq G$: $H\neq\emptyset$ and closed under the operation $\Rightarrow$ $H\le G$.

\subsect{Important Subgroups}
\begin{itemize}[nosep,leftmargin=*]
\item $\langle a\rangle = \{a^n : n\in\mathbb{Z}\}$: cyclic subgroup generated by $a$.
\item Center: $Z(G) = \{g\in G : gx=xg\;\forall x\in G\}$.
\item Centralizer: $C_G(a)=\{g\in G: ga=ag\}$.
\item $Z(G)\le C_G(a)\le G$ for every $a\in G$.
\item $Z(G)$ is always a \textbf{normal} subgroup of $G$.
\end{itemize}

%======================================================================
\sect{3. Cyclic Groups}
%======================================================================

$G=\langle a\rangle$: every element is a power of $a$.
\begin{itemize}[nosep,leftmargin=*]
\item Cyclic $\Rightarrow$ abelian (converse is FALSE: $V_4$).
\item $|\mathbb{Z}/n\mathbb{Z}|=n$. Generators: $k$ with $\gcd(k,n)=1$.
\item Number of generators of $\mathbb{Z}/n\mathbb{Z}$ is $\varphi(n)$ (Euler's totient).
\item Every subgroup of a cyclic group is cyclic.
\item $\mathbb{Z}/n\mathbb{Z}$ has exactly one subgroup of order $d$ for each $d\mid n$.
\item $\mathbb{Z}/m\mathbb{Z}\times\mathbb{Z}/n\mathbb{Z}\cong\mathbb{Z}/mn\mathbb{Z}$ $\iff$ $\gcd(m,n)=1$.
\end{itemize}

\subsect{Classification}
Every cyclic group is isomorphic to $\mathbb{Z}$ (infinite) or $\mathbb{Z}/n\mathbb{Z}$ (finite, order $n$).

%======================================================================
\sect{4. Cosets \& Lagrange's Theorem}
%======================================================================

\subsect{Cosets}
For $H\le G$ and $a\in G$:
\begin{itemize}[nosep,leftmargin=*]
\item Left coset: $aH = \{ah : h\in H\}$.
\item Right coset: $Ha = \{ha : h\in H\}$.
\item All left cosets have the same size: $|aH|=|H|$.
\item Left cosets partition $G$.
\end{itemize}

\subsect{Lagrange's Theorem}
If $H\le G$ and $|G|<\infty$, then $|H|$ divides $|G|$.\\
Index: $[G:H]=|G|/|H|=$ number of distinct cosets.

\textbf{Consequences:}
\begin{itemize}[nosep,leftmargin=*]
\item $\operatorname{ord}(a)$ divides $|G|$.
\item $a^{|G|}=e$ for all $a\in G$.
\item A group of prime order $p$ is cyclic ($\cong\mathbb{Z}/p\mathbb{Z}$).
\item If $|G|=p$ (prime), the only subgroups are $\{e\}$ and $G$.
\item \textbf{Converse is FALSE:} $A_4$ has order 12 but no subgroup of order 6.
\end{itemize}

%======================================================================
\sect{5. Normal Subgroups \& Quotient Groups}
%======================================================================

\subsect{Normal Subgroup}
$N\trianglelefteq G$ if $gNg^{-1}=N$ for all $g\in G$, equivalently $gN=Ng$ for all $g$.
\begin{itemize}[nosep,leftmargin=*]
\item Every subgroup of an abelian group is normal.
\item $Z(G)\trianglelefteq G$ always.
\item $\ker\phi\trianglelefteq G$ for any homomorphism $\phi$.
\item Index-2 subgroups are always normal.
\item $A_n\trianglelefteq S_n$ with $[S_n:A_n]=2$.
\end{itemize}

\subsect{Quotient Group}
$G/N = \{gN : g\in G\}$ with operation $(aN)(bN)=(ab)N$.\\
$|G/N|=[G:N]=|G|/|N|$.

\subsect{Simple Groups}
$G$ is simple if its only normal subgroups are $\{e\}$ and $G$.
\begin{itemize}[nosep,leftmargin=*]
\item $\mathbb{Z}/p\mathbb{Z}$ is simple for $p$ prime.
\item $A_n$ is simple for $n\ge 5$.
\end{itemize}

%======================================================================
\sect{6. Homomorphisms \& Isomorphisms}
%======================================================================

\subsect{Homomorphism}
$\phi\colon G\to H$ with $\phi(ab)=\phi(a)\phi(b)$ for all $a,b$.
\begin{itemize}[nosep,leftmargin=*]
\item $\phi(e_G)=e_H$.
\item $\phi(a^{-1})=\phi(a)^{-1}$.
\item $\operatorname{ord}(\phi(a))$ divides $\operatorname{ord}(a)$.
\item $\ker\phi=\{g\in G:\phi(g)=e_H\}$ is a normal subgroup of $G$.
\item $\operatorname{im}\phi=\phi(G)$ is a subgroup of $H$.
\item $\phi$ injective $\iff$ $\ker\phi=\{e\}$.
\end{itemize}

\subsect{First Isomorphism Theorem}
$G/\ker\phi \cong \operatorname{im}\phi$.

\subsect{Counting}
$|\operatorname{im}\phi|$ divides both $|G|$ and $|H|$.\\
$|G|=|\ker\phi|\cdot|\operatorname{im}\phi|$.

%======================================================================
\sect{7. Permutation Groups}
%======================================================================

\subsect{Symmetric Group $S_n$}
$|S_n|=n!$. Elements are bijections $\{1,\ldots,n\}\to\{1,\ldots,n\}$.

\subsect{Cycle Notation}
$(1\;2\;3)$ means $1\mapsto 2\mapsto 3\mapsto 1$.\\
Disjoint cycles commute. Every permutation is a product of disjoint cycles.

\subsect{Key Facts}
\begin{itemize}[nosep,leftmargin=*]
\item Order of a permutation $=$ lcm of the lengths of its disjoint cycles.
\item A $k$-cycle is a product of $k-1$ transpositions.
\item $\sigma$ is even if it's a product of an even number of transpositions.
\item $\operatorname{sign}(\sigma)=(-1)^{n-c}$ where $c=$ number of cycles (including 1-cycles).
\end{itemize}

\subsect{Alternating Group $A_n$}
$A_n = \{\sigma\in S_n : \sigma\text{ is even}\}$, $\;|A_n|=n!/2$.
\begin{itemize}[nosep,leftmargin=*]
\item $A_3\cong\mathbb{Z}/3\mathbb{Z}$ (cyclic).
\item $A_4$: order 12, has no subgroup of order 6.
\item $A_5$: simple, order 60.
\end{itemize}

%======================================================================
\sect{8. Group Actions \& Sylow Theorems}
%======================================================================

\subsect{Class Equation}
$|G| = |Z(G)| + \sum [G:C_G(a_i)]$\\
where the sum is over one representative $a_i$ from each non-trivial conjugacy class.

\subsect{$p$-Groups}
$|G|=p^k$ for some prime $p$.
\begin{itemize}[nosep,leftmargin=*]
\item $p$-groups have non-trivial center: $Z(G)\neq\{e\}$.
\item Groups of order $p^2$ are abelian (either $\mathbb{Z}/p^2\mathbb{Z}$ or $(\mathbb{Z}/p\mathbb{Z})^2$).
\end{itemize}

\subsect{Sylow Theorems}
Let $|G|=p^a m$ with $\gcd(p,m)=1$.
\begin{enumerate}[nosep,leftmargin=*]
\item There exists a subgroup of order $p^a$ (Sylow $p$-subgroup).
\item All Sylow $p$-subgroups are conjugate.
\item $n_p$ (number of Sylow $p$-subgroups) satisfies $n_p\mid m$ and $n_p\equiv 1\pmod{p}$.
\end{enumerate}
\textbf{Key application:} If $n_p=1$, the Sylow $p$-subgroup is \textbf{normal}.

\subsect{Quick Examples}
\begin{itemize}[nosep,leftmargin=*]
\item $|G|=15=3\cdot 5$: $n_3\mid 5,\; n_3\equiv 1\pmod{3}\Rightarrow n_3=1$; $n_5\mid 3,\; n_5\equiv 1\pmod{5}\Rightarrow n_5=1$. So $G\cong\mathbb{Z}/15\mathbb{Z}$.
\item $|G|=pq$ ($p<q$ primes, $p\nmid q-1$): $G$ is cyclic.
\end{itemize}

%======================================================================
\sect{9. Rings}
%======================================================================

\subsect{Ring Axioms}
$(R,+,\cdot)$: $(R,+)$ is an abelian group, $\cdot$ is associative, and distribution holds.\\
If $\cdot$ is commutative: \textbf{commutative ring}. If $\exists\,1$: \textbf{ring with unity}.

\subsect{Common Examples}
\begin{itemize}[nosep,leftmargin=*]
\item $\mathbb{Z}$, $\mathbb{Q}$, $\mathbb{R}$, $\mathbb{C}$: commutative rings with unity.
\item $\mathbb{Z}/n\mathbb{Z}$: commutative ring with unity.
\item $M_n(\mathbb{R})$: ring with unity, NOT commutative for $n\ge 2$.
\item $\mathbb{Z}[x]$: polynomial ring, commutative with unity.
\item $2\mathbb{Z}$: commutative ring \textbf{without} unity.
\end{itemize}

\subsect{Special Elements}
\begin{itemize}[nosep,leftmargin=*]
\item \textbf{Unit:} $a\in R$ with $ab=ba=1$ for some $b$.
\item \textbf{Zero divisor:} $a\neq 0$ with $ab=0$ for some $b\neq 0$.
\item \textbf{Nilpotent:} $a^n=0$ for some $n\ge 1$.
\item A unit is never a zero divisor.
\end{itemize}

\subsect{Integral Domain}
Commutative ring with unity, \textbf{no zero divisors}.\\
$\mathbb{Z}$, $\mathbb{Z}[x]$, $\mathbb{Z}/p\mathbb{Z}$ ($p$ prime) are integral domains.\\
$\mathbb{Z}/n\mathbb{Z}$ ($n$ composite) is NOT an integral domain.

\subsect{Ideals}
$I\subseteq R$ is an ideal if: $I$ is a subgroup of $(R,+)$ and $ra,ar\in I$ for all $r\in R,\;a\in I$.
\begin{itemize}[nosep,leftmargin=*]
\item $(n) = n\mathbb{Z}$: ideal of $\mathbb{Z}$ generated by $n$.
\item $R/I$ is a ring (quotient ring).
\item $\ker\phi$ is an ideal for any ring homomorphism $\phi$.
\end{itemize}

\subsect{Types of Ideals}
\begin{itemize}[nosep,leftmargin=*]
\item \textbf{Principal:} $I=(a)=\{ra:r\in R\}$.
\item \textbf{Maximal:} $I\neq R$ and no ideal strictly between $I$ and $R$.
\item \textbf{Prime:} $ab\in I\Rightarrow a\in I$ or $b\in I$.
\item Maximal $\Rightarrow$ prime (converse is FALSE in general).
\item $R/I$ is a field $\iff$ $I$ is maximal.
\item $R/I$ is an integral domain $\iff$ $I$ is prime.
\end{itemize}

\subsect{Special Rings}
\begin{itemize}[nosep,leftmargin=*]
\item \textbf{PID} (principal ideal domain): every ideal is principal.\\
  $\mathbb{Z}$ and $F[x]$ ($F$ a field) are PIDs. $\;\mathbb{Z}[x]$ is NOT a PID.
\item \textbf{UFD} (unique factorization domain): every non-zero non-unit factors uniquely into irreducibles.\\
  PID $\Rightarrow$ UFD (converse is FALSE: $\mathbb{Z}[x]$).
\item \textbf{ED} (Euclidean domain): has a division algorithm.\\
  ED $\Rightarrow$ PID $\Rightarrow$ UFD $\Rightarrow$ integral domain.
\end{itemize}

%======================================================================
\sect{10. Fields}
%======================================================================

\subsect{Definition}
A commutative ring with unity in which every non-zero element is a unit.\\
Equivalently: exactly two ideals, $(0)$ and $R$.

\subsect{Characteristic}
$\operatorname{char}(F)=$ smallest $n>0$ s.t.\ $n\cdot 1=0$, or $0$ if no such $n$ exists.\\
$\operatorname{char}(F)$ is always $0$ or prime.
\begin{itemize}[nosep,leftmargin=*]
\item $\operatorname{char}(\mathbb{Q})=\operatorname{char}(\mathbb{R})=\operatorname{char}(\mathbb{C})=0$.
\item $\operatorname{char}(\mathbb{Z}/p\mathbb{Z})=p$.
\item In characteristic $p$: $(a+b)^p = a^p+b^p$ (Frobenius endomorphism).
\end{itemize}

\subsect{Finite Fields}
\begin{itemize}[nosep,leftmargin=*]
\item A finite field has $p^n$ elements for some prime $p$ and $n\ge 1$.
\item For each prime power $q=p^n$, there exists a \textbf{unique} (up to isomorphism) field $\mathbb{F}_q$ of order $q$.
\item $\mathbb{F}_q^*$ (multiplicative group) is \textbf{cyclic} of order $q-1$.
\item $\mathbb{F}_p\cong\mathbb{Z}/p\mathbb{Z}$.
\end{itemize}

\subsect{Field Extensions (Basics)}
$F\subseteq K$: $K$ is an extension of $F$.\\
$[K:F]=\dim_F K$ (degree of extension).
\begin{itemize}[nosep,leftmargin=*]
\item $[\mathbb{C}:\mathbb{R}]=2$, $\;[\mathbb{R}:\mathbb{Q}]=\infty$.
\item $[\mathbb{Q}(\sqrt{2}):\mathbb{Q}]=2$; minimal polynomial of $\sqrt{2}$ is $x^2-2$.
\item Tower law: $[K:F]=[K:E]\cdot[E:F]$.
\item If $[K:F]$ is prime, there is no intermediate field.
\end{itemize}

\subsect{Irreducibility Tests}
\begin{itemize}[nosep,leftmargin=*]
\item Degree 2 or 3: irreducible over $F$ $\iff$ no root in $F$.
\item \textbf{Eisenstein:} $f(x)=a_nx^n+\cdots+a_0\in\mathbb{Z}[x]$. If $\exists$ prime $p$ with $p\nmid a_n$, $p\mid a_i$ for $i<n$, and $p^2\nmid a_0$, then $f$ is irreducible over $\mathbb{Q}$.
\item Reduction mod $p$: If $f\bmod p$ is irreducible in $\mathbb{F}_p[x]$ (and same degree), then $f$ is irreducible over $\mathbb{Q}$.
\end{itemize}

%======================================================================
\sect{11. Key Classification Results}
%======================================================================

\subsect{Groups of Small Order}
\begin{tabular}{@{}cl@{}}
\toprule
$|G|$ & Groups (up to isomorphism) \\
\midrule
1 & $\{e\}$ \\
2 & $\mathbb{Z}/2\mathbb{Z}$ \\
3 & $\mathbb{Z}/3\mathbb{Z}$ \\
4 & $\mathbb{Z}/4\mathbb{Z}$, \;$V_4$ \\
5 & $\mathbb{Z}/5\mathbb{Z}$ \\
6 & $\mathbb{Z}/6\mathbb{Z}$, \;$S_3$ \\
7 & $\mathbb{Z}/7\mathbb{Z}$ \\
8 & $\mathbb{Z}/8\mathbb{Z}$, \;$\mathbb{Z}/4\mathbb{Z}\times\mathbb{Z}/2\mathbb{Z}$, \;$(\mathbb{Z}/2\mathbb{Z})^3$, \;$D_4$, \;$Q_8$ \\
\bottomrule
\end{tabular}

\subsect{Fundamental Theorem of Finite Abelian Groups}
Every finite abelian group is isomorphic to a direct product of cyclic groups of prime-power order:
\[
G\cong \mathbb{Z}/p_1^{a_1}\mathbb{Z}\times\cdots\times\mathbb{Z}/p_k^{a_k}\mathbb{Z}.
\]
This decomposition is unique (up to reordering).

\textbf{Example:} Abelian groups of order 12:
$\mathbb{Z}/12\mathbb{Z}$ \;and\; $\mathbb{Z}/2\mathbb{Z}\times\mathbb{Z}/6\mathbb{Z}$ $(\cong\mathbb{Z}/2\mathbb{Z}\times\mathbb{Z}/2\mathbb{Z}\times\mathbb{Z}/3\mathbb{Z})$.

\bigskip
\hrule
\smallskip
\begin{center}
\textit{Hierarchy:}\; Field $\subset$ ED $\subset$ PID $\subset$ UFD $\subset$ Integral Domain $\subset$ Commutative Ring.
\end{center}

\end{document}
